\section{Exact stationary endpoint inversion by flight-age integration}
\label{sec:exact-stationary-endpoint}
The return record from the preceding section and the observation at a
fixed physical time are different microscopic events. In this section
we invert the latter directly. Integrating the actual initial flight
age produces an overlap of two intervals. This overlap has two
distributional derivatives with a uniform measure bound, independently
of the number of intervening collisions. The initial and terminal
cell corrections are retained exactly, not bounded and discarded.

Throughout this section $m\ge1$ is a prescribed \emph{physical collision
count}, not a return count. Put
\[
 S_{m,R}=\sum_{j=0}^{m-1}\tau_R\circ T_R^j,\qquad
 K^{\rm c}_{m,R}=\sum_{j=0}^{m-1}\kappa_R\circ T_R^j,
 \qquad 0<\tau_-\le\tau_R\le\tau_+.
\]
Here $\kappa_R$ is the physical next-disk label and $\tau_-,\tau_+$
are uniform finite-horizon bounds. The equilibrium probability in
current-collision coordinates is, by
Lemma~\ref{lem:stationary-length-bias},
\begin{equation}\label{eq:stationary-endpoint-source}
 d\mathbb P_R(x,a)=\bar\tau_R^{-1}d\nu(x)\,da,
                    \qquad 0\le a<\tau_R(x).
\end{equation}
This is the same stationary trajectory as in the return-suspension
coordinates. No new initial point is sampled.

\subsection{Cell labels during a physical flight}
Fix the fundamental-cell convention used for $W_R(t)$. We take a
bounded convex polygonal fundamental cell $\mathcal P$; in particular
one may use a fixed half-open lattice parallelogram. Let
$\mathcal L:\R^2\to\R^2$ send the standard integer basis to the
triangular lattice basis. For the lifted flight beginning at the disk
with center zero define
\begin{equation}\label{eq:within-flight-cell}
 d_R(x,a)=d\quad\Longleftrightarrow\quad
 Rn_\alpha+a v(x)\in\mathcal P+\mathcal Ld,
                \qquad d\in\Z^2,
\end{equation}
with the fixed half-open boundary convention. The same cell convention
is used at the two observation times.

\begin{lemma}[A uniform finite flight partition]
\label{lem:finite-flight-cell-partition}
There is a fixed finite set $\mathcal D\subset\Z^2$ such that, for
almost every regular $x$ and every $R$, the sets
\[
 I_{R,d}(x)=\{a\in[0,\tau_R(x)):d_R(x,a)=d\},\qquad d\in\mathcal D,
\]
are intervals or empty and partition the flight up to endpoints.
Write $J_* = |\mathcal D|$. Both $J_*$ and $\max_{d\in\mathcal D}|d|$
are uniform in $R$.

If $C_R(t)=m$ for a trajectory starting at $(x,a)$, its terminal
collision state is $y=T_R^mx$, its terminal age is
$v=t+a-S_{m,R}(x)$, and exactly
\begin{equation}\label{eq:exact-physical-cell-correction}
 0\le v<\tau_R(y),\qquad
 W_R(t)=K^{\rm c}_{m,R}(x)+d_R(y,v)-d_R(x,a).
\end{equation}
\end{lemma}
\begin{proof}
The segment in \eqref{eq:within-flight-cell} is contained in the fixed
ball of radius $47/100+\tau_+$. Only finitely many translates of the
bounded fundamental cell meet that ball; take their labels as
$\mathcal D$. Convexity makes the intersection of each cell with a
straight flight an interval. Flights lying along a cell edge form a
null set and may be assigned the fixed boundary convention. This
proves the partition and uniformity without differentiating a cell
crossing time.

The $j$th future collision occurs at $S_{j,R}(x)-a$. These times are
strictly increasing, so the terminal age condition is equivalent to
$C_R(t)=m$, apart from suspension-null endpoints. The terminal lifted
position is
$\mathcal L K^{\rm c}_{m,R}(x)+Rn_{\alpha_y}+v v(y)$.
Subtracting the initial cell label from its cell label gives
\eqref{eq:exact-physical-cell-correction}. In particular the last two
terms cannot be omitted at a singleton lattice label.
\end{proof}

For intervals $I,J\subset\R$ put
\begin{equation}\label{eq:flight-overlap-kernel}
 H_{I,J}(s)=\int_\R\one_I(a)\one_J(s+a)\,da
                         =(\one_{-I}*\one_J)(s).
\end{equation}
The value does not depend on choices at interval endpoints. Define
$p_{m,R}(k,t)=\mathbb P_R\{W_R(t)=k,C_R(t)=m\}$ for $t\ge0$ and
extend it by zero to $t<0$.

\begin{theorem}[Exact overlap formula and uniform physical-time regularity]
\label{thm:stationary-overlap-regularity}
For all $m\ge1$ the profiles $p_{m,R}(k,\cdot)$ have continuous,
compactly supported representatives given by
\begin{align}
 p_{m,R}(k,t)
 =\frac1{\bar\tau_R}\int_M\sum_{d,e\in\mathcal D}
 &\one_{\{K^{\rm c}_{m,R}(x)+e-d=k\}}\notag\\[-1mm]
 &\cdot H_{I_{R,d}(x),I_{R,e}(T_R^mx)}(t-S_{m,R}(x))\,d\nu(x).
 \label{eq:exact-overlap-probability}
\end{align}
There are constants $A_0,A_1,A_2,L_0<\infty$, independent of $m,R$,
for which
\begin{align}
 \sum_k\|p_{m,R}(k,\cdot)\|_1&\le A_0,&
 \sum_k\|\partial_t p_{m,R}(k,\cdot)\|_1&\le A_1,
                                      \label{eq:stationary-time-W1}\\
 \sum_k\|D_t^2p_{m,R}(k,\cdot)\|_{\TV}&\le A_2,&
 \sum_k|p_{m,R}(k,t)-p_{m,R}(k,s)|&\le L_0|t-s|.
                                      \label{eq:stationary-time-curvature}
\end{align}
One may take, writing $\bar\tau_- =\min_R\bar\tau_R$,
\[
 A_0=\tau_+^2/\bar\tau_-,\quad
 A_1=2J_*\tau_+/\bar\tau_-,\quad
 A_2=4J_*^2/\bar\tau_-,\quad L_0=2J_*/\bar\tau_-.
\]
Here $D_t^2p$ is a finite signed measure; an $L^1$ density for this
second derivative is neither needed nor asserted. The profiles are
probabilities at fixed observation times, not a probability density
of total mass one on $\Z^2\times\R$.
\end{theorem}
\begin{proof}
Insert \eqref{eq:exact-physical-cell-correction} into
\eqref{eq:stationary-endpoint-source} and integrate the initial age.
This is exactly \eqref{eq:exact-overlap-probability}. Since
$S_m\ge\tau_R(x)$, its integrand vanishes for $t<0$; it is supported
in $0\le t\le(m+1)\tau_+$. There are only finitely many lattice
labels for each $m$, by finite horizon.

For an interval, $D\one_I$ is the difference of its endpoint masses,
with total variation at most two. Consequently
\[
 \|H_{I,J}\|_1=|I||J|,\qquad
 \|H'_{I,J}\|_1\le2|I|,\qquad
 \|D^2H_{I,J}\|_{\TV}\le4.
\]
Apply these inequalities before integrating over $x$. Summing over
$k$ counts each pair $(d,e)$ exactly once. The interval partitions give
$\sum_d|I_{R,d}(x)|=\tau_R(x)$ and the analogous terminal identity.
They yield all three bounds $A_0,A_1,A_2$. Differentiation is justified
in the distributional sense by testing against compactly supported
smooth functions and using these finite absolute bounds.

For the Lipschitz assertion it is sharper to work on the source.
Denote the integrand event indicator at time $t$ by $a_k(t;x,a)$.
The unique initial interval and the final partition imply
\begin{align*}
 \sum_k|a_k(t;x,a)-a_k(s;x,a)|
 \le \sum_{d,e}\one_{I_{R,d}(x)}(a)
 \big|\one_{I_{R,e}(T^mx)}(t+a-S_m)
       -\one_{I_{R,e}(T^mx)}(s+a-S_m)\big|.
\end{align*}
Sum first over $d$, integrate in $a$, and enlarge its domain to the
real line. Translation of each final interval costs at most $2|t-s|$.
Integrating over $\nu$ proves the last bound with $L_0$. It also proves
continuity of the formula, including at $t=0$.
\end{proof}

\subsection{An absolutely convergent physical inversion}
Use $\widehat f(b)=\int e^{ibt}f(t)\,dt$. For $u\in\T^2$ and
$b\in\R$ define the single-flight amplitudes
\begin{align}
 A^-_R(x;u,b)&=\int_0^{\tau_R(x)}
           e^{-i(u\cdot d_R(x,a)+ba)}\,da,\notag\\
 A^+_R(x;u,b)&=\int_0^{\tau_R(x)}
           e^{ i(u\cdot d_R(x,a)+ba)}\,da,
                                    \label{eq:single-flight-amplitudes}\\
 \mathcal F_{m,R}(u,b)&=\int_M
 e^{i(u\cdot K^{\rm c}_{m,R}+bS_{m,R})}
 A^-_R(x;u,b)A^+_R(T_R^mx;u,b)\,d\nu(x).
                                    \label{eq:stationary-physical-transform}
\end{align}
The torus is represented by $[-\pi,\pi]^2$. For real frequencies the
two single-flight amplitudes are complex conjugates at the same state.

\begin{theorem}[Uniform pointwise far-frequency bound]
\label{thm:stationary-polynomial-band}
The probability of the original singleton physical event satisfies
\begin{equation}\label{eq:stationary-absolute-inversion}
 p_{m,R}(k,t)=\frac1{(2\pi)^3\bar\tau_R}
 \int_{[-\pi,\pi]^2}\int_\R
 e^{-i(u\cdot k+bt)}\mathcal F_{m,R}(u,b)\,db\,du.
\end{equation}
The integral is absolutely convergent. If $p^{\rm sharp}_{m,R,B}$
denotes the same inverse restricted to $|b|\le B$, then
\begin{equation}\label{eq:stationary-sharp-pointwise-tail}
 \sum_k\|p_{m,R}(k,\cdot)-p^{\rm sharp}_{m,R,B}(k,\cdot)\|_\infty
                          \le A_2/(\pi B),\qquad B>0.
\end{equation}
In particular $B_m=m^{P+3/2}$ gives error $O(m^{-P})$ after
multiplication by the physical local scale $m^{3/2}$.
\end{theorem}
\begin{proof}
Each amplitude is a sum of at most $J_*$ interval integrals, so
\begin{equation}\label{eq:flight-amplitude-decay}
 |A_R^\pm(x;u,b)|\le\min\{\tau_+,2J_*/|b|\}.
\end{equation}
The product in \eqref{eq:stationary-physical-transform} is therefore
bounded by $\min\{\tau_+^2,4J_*^2/b^2\}$, uniformly in $m,R,u$.
Fourier transformation of \eqref{eq:flight-overlap-kernel} shows that
$\mathcal F_{m,R}/\bar\tau_R$ is exactly the Fourier transform of the
finite family of time profiles. Fourier inversion and torus
orthogonality prove \eqref{eq:stationary-absolute-inversion}; continuity
from Theorem~\ref{thm:stationary-overlap-regularity} makes the identity
pointwise. No dynamical decay inside the finite band is invoked.

For the stronger summed bound, integration by parts in the sense of
distributions gives, for $b\ne0$,
\[
 \sum_k|\widehat p_{m,R}(k,b)|
 \le |b|^{-2}\sum_k\|D_t^2p_{m,R}(k,\cdot)\|_{\TV}
 \le A_2/b^2.
\]
Integrating over $|b|>B$ with the one-dimensional inverse factor
$1/(2\pi)$ proves \eqref{eq:stationary-sharp-pointwise-tail}.
The factor from the lattice inversion has already been accounted for
in each Fourier coefficient. This avoids a factor equal to the number
of spatial labels.
\end{proof}

\subsection{A positive likelihood on the same physical trajectory}
Let $k_B$ and $\eta$ be the positive kernel and its transform from
Lemma~\ref{lem:positive-kernel-moments}. For a set
$\mathcal K\subset\Z^2$ define, on the original stationary source,
\begin{equation}\label{eq:physical-time-positive-likelihood}
 A_{m,\mathcal K}(t)=\{C_R(t)=m,W_R(t)\in\mathcal K\},\qquad
 \lambda_{m,\mathcal K,B}(t;\omega)
   =\int_\R k_B(h)\one_{A_{m,\mathcal K}(t-h)}(\omega)\,dh.
\end{equation}
Negative-time events in this formula are zero. These are numerical
weights on one trajectory. They do not change its physical dynamics,
its age, its lattice convention or the target event.

\begin{theorem}[Polynomial-band reconstruction of the original event]
\label{thm:physical-time-selector-reconstruction}
For every $m\ge1$, $R$, $t\ge0$, $B>0$ and spatial set $\mathcal K$,
\begin{equation}\label{eq:physical-time-source-TV}
 0\le\lambda_{m,\mathcal K,B}\le1,\qquad
 \|\one_{A_{m,\mathcal K}(t)}-\lambda_{m,\mathcal K,B}(t)\|_{L^1(\mathbb P_R)}
                          \le L_0\kappa_1/B=:E_B^{\rm phys}.
\end{equation}
Consequently, for every bounded measurable trajectory statistic $W$,
\begin{equation}\label{eq:physical-time-all-selectors}
 |\mathbb E_R[W\one_{A_{m,\mathcal K}(t)}]
     -\mathbb E_R[W\lambda_{m,\mathcal K,B}(t)]|
                           \le\|W\|_\infty E_B^{\rm phys}.
\end{equation}
Uniformity permits $W$ to depend on $m,t,B$. It follows by source
contraction, not a spectral estimate for the weighted transform.

For $\mathcal K=\{k\}$ the second expectation has the finite formula
\begin{equation}\label{eq:physical-time-weighted-formula}
 \begin{aligned}
 \frac1{(2\pi)^3\bar\tau_R}
 \int_{[-\pi,\pi]^2}\int_{-B}^B
 &e^{-i(u\cdot k+bt)}\eta(b/B)\\
 {}\cdot\int_M &e^{i(u\cdot K^{\rm c}_{m,R}+bS_{m,R})}
 A^-_{R,W}(x;u,b)A^+_R(T_R^mx;u,b)\,d\nu(x)\,db\,du,
 \end{aligned}
\end{equation}
where
\[
 A^-_{R,W}(x;u,b)=\int_0^{\tau_R(x)}
             W(x,a)e^{-i(u\cdot d_R(x,a)+ba)}\,da.
\]
The statistic $W$ is held fixed while taking the observation-time
transform. A fixed bounded target-time interval $I$ may also be
integrated: its absolute weighted error is at most
$|I|\|W\|_\infty E_B^{\rm phys}$.
\end{theorem}
\begin{proof}
The source argument in the proof of
Theorem~\ref{thm:stationary-overlap-regularity}, before averaging the
event indicators, shows the stronger estimate
\[
 \sum_k\mathbb E_R|\one_{A_{m,\{k\}}(t)}-
                       \one_{A_{m,\{k\}}(s)}|\le L_0|t-s|
\]
for all real $t,s$, with the zero extension already established.
A spatial set can only reduce the corresponding bound. Average it
against $k_B(t-s)$ and use positivity and
$\int |h|k_B(h)\,dh=\kappa_1/B$. This proves
\eqref{eq:physical-time-source-TV}. Multiplication by $W$ and then
integration prove \eqref{eq:physical-time-all-selectors}. In particular
this is total variation of unnormalized measures on the original
stationary source, not merely a difference of scalar probabilities.

For fixed $(x,a)$ and $m$, the time event is supported in an interval
of length at most $\tau_+$. Its time Fourier transform is obtained by
the substitution $t=S_m+v-a$. Integrate $v$ over the final cell
intervals and $a$ over the initial ones. The factor $W(x,a)$ remains
inside the latter integral and gives $A^-_{R,W}$. Invert the positive
kernel and the two lattice coordinates to obtain
\eqref{eq:physical-time-weighted-formula}. The band is finite and all
age integrals and the statistic are bounded, so Fubini is legitimate.
Finally integrate the established error over $I$.
\end{proof}

\begin{corollary}[Denominator-sensitive physical posteriors]
\label{cor:direct-physical-posterior}
Let $A=A_{m,\mathcal K}(t)$, $p_A=\mathbb P_R(A)$ and
$p_{A,B}=\mathbb E_R\lambda_{m,\mathcal K,B}(t)$. If
$p_{A,B}>E_B^{\rm phys}$, then
\[
 p_A\ge p_{A,B}-E_B^{\rm phys}>0,\qquad
 \left\|\mathbb P_R(\cdot\mid A)-
          \frac{\lambda_{m,\mathcal K,B}(t)}{p_{A,B}}\mathbb P_R\right\|_{\TV}
                \le\frac{2E_B^{\rm phys}}{p_{A,B}-E_B^{\rm phys}}.
\]
If instead a lower bound for $p_A>E_B^{\rm phys}$ is known, the right
side may be replaced by $2E_B^{\rm phys}/p_A$. These assertions hold
also after multiplication by any fixed nonnegative selector at most
one, using its own denominator. For $p_A\ge c m^{-3/2}$ and
$B_m=m^{P+3/2}$, the relative error is $O(m^{-P})$.
\end{corollary}
\begin{proof}
Apply the normalization inequality in
Corollary~\ref{cor:positive-microscopic-posterior} to
\eqref{eq:physical-time-source-TV}. Multiplication by a selector
contracts the source $L^1$ error. The last assertion is division by
the explicitly stated denominator bound, not a proof of that bound.
\end{proof}

\paragraph{Which tail has been closed.}
Equations~\eqref{eq:stationary-sharp-pointwise-tail} and
\eqref{eq:physical-time-source-TV} use a polynomial bandwidth for the
\emph{original stationary observation} at a fixed collision count.
They require no removal of a small-mass source set, no long-itinerary
coarea derivative bound, and no comparison with an event at the preceding
section origin. The time profiles have uniform distributional
curvature because two physical flight ages have been integrated.
The return density in Theorem W has a different output and does not
inherit this estimate. In particular a bound for these profiles is not
a pointwise estimate for the four-coordinate return density.
